% changelog: "1.0.0, 2024-05-08, Approvazione per RTB"

\documentclass[8pt]{article}
\usepackage[italian]{babel}
\usepackage[utf8]{inputenc}
\usepackage[letterpaper, left=1in, right=1in, bottom=0.75in, top=0.75in]{geometry}
\usepackage{amsmath}
\usepackage{subfiles}
\usepackage{lipsum}
\usepackage{csquotes}
\usepackage{amsfonts}
\usepackage[sfdefault]{plex-sans}
\usepackage{float}
\usepackage{pifont}
\usepackage{mathabx}
\usepackage[euler]{textgreek}
\usepackage{makecell}
\usepackage{tikz}
\usepackage{wrapfig}
\usepackage{siunitx}
\usepackage{amssymb} 
\usepackage{tabularx}
\usepackage{adjustbox}
\usepackage[document]{ragged2e}
\usepackage{floatflt}
\usepackage[hidelinks]{hyperref}
\usepackage{graphicx}
\usepackage{hyperref}
\setcounter{tocdepth}{4}
\usepackage{caption}
\usepackage{multicol}
\usepackage{tikz}
\setlength\parindent{0pt}
\captionsetup{font=footnotesize}
\usepackage{fancyhdr} 
\usepackage{graphicx}
\usepackage{capt-of}% 
\usepackage{booktabs}
\usepackage{varwidth}


\newcommand{\customtitle}{VERBALE ESTERNO DEL 2024-05-03}

\fancypagestyle{mystyle}{
	\fancyhf{} 
	\fancyhead[R]{\includegraphics[height=1cm]{../../../template/images/logos/NaN1fy_logo.png}} 
	\fancyhead[L]{\leftmark} 
	\renewcommand{\headrulewidth}{1pt} 
	\fancyhead[L]{\customtitle} 
	\renewcommand{\headsep}{1.3cm} 
	\fancyfoot[C]{\thepage} 
}

\newcommand{\signatureline}[1]{%
	 \par\vspace{0.5cm}
	\noindent\makebox[\linewidth][r]{\rule{0.2\textwidth}{0.5pt}\hspace{3cm}\makebox[0pt][r]{\vspace{3pt}\footnotesize #1}}%
}

\newcommand{\glossterm}[1]{#1\textsuperscript{G}} % inserisci \glossterm{termine}

\begin{document}
\definecolor{myblue}{RGB}{23,103,162}
\begin{titlepage}
	\begin{tikzpicture}[remember picture, overlay]
      \node[anchor=south east, opacity=0.2, yshift = -4cm, xshift= 2em] at (current page.south east) {\includegraphics[width=0.7\textwidth, trim=0cm 0cm 5cm 0cm, clip]{../../../template/images/logos/Universita_Padova_transparent.png}};
		\node[anchor=north west, opacity=1, yshift = 4.2cm, xshift= 1.4cm, scale=1.6] at (current page.south west) {\includegraphics[width=4cm]{../../../template/images/logos/NaN1fy_logo.png}};
	\end{tikzpicture}
	
	\begin{minipage}[t]{0.47\textwidth}
		{\large{\textsc{Destinatari}}
			\vspace{3mm}
			\\ \large{\textsc{Prof. Tullio Vardanega}}
			\\ \large{\textsc{Prof. Riccardo Cardin}}
		}
	\end{minipage}
	\hfill
	\begin{minipage}[t]{0.47\textwidth}\raggedleft
		{\large{\textsc{Redattori}}
			\vspace{3mm}
			{\\\large{\textsc{Pietro Busato}\\}} % massimo due 
		}
		\vspace{8mm}
		
		{\large{\textsc{Verificatori}}
			\vspace{3mm}
			{\\\large{\textsc{Guglielmo Barison}\\}} % massimo due 
			{\large{\textsc{Veronica Tecchiati}}}
			
		}
		\vspace{4mm}\vspace{4mm}
	\end{minipage}
	\vspace{4cm}
	\begin{center}
		\begin{flushright}
			{\fontsize{30pt}{52pt}\selectfont \textbf{Verbale Esterno Del\\2024-05-03\\}} % o ESTERNO
		\end{flushright}
		\vspace{3cm}
	\end{center}
	\vspace{8.5 cm}
	{\small \textsc{\href{mailto: nan1fyteam.unipd@gmail.com}{nan1fyteam.unipd@gmail.com}}}
\end{titlepage}

\pagestyle{mystyle}
\section*{Registro delle Modifiche}
\begin{table}[ht!]	
	\centering
	\begin{tabular}{p{1.2cm} p{2cm} p{6cm} p{3cm} p{2cm}}
		\toprule
		\textbf{Versione}& \textbf{Data} & \textbf{Descrizione} & \textbf{Autore} & \textbf{Ruolo} \\
		\midrule	 
    1.0.0 & 2024-05-08 & \textbf{Approvazione per RTB}  &  & \\\\
    0.0.1 & 2024-05-08 & Verifica completa, correzioni minori.  & Guglielmo Barison, Veronica Tecchiati & Verificatore \\\\
    0.0.0 & 2024-05-07 & Stesura documento.  & Pietro Busato & Redattore \\
		\bottomrule
		% Ruolo Redattore o Verificatore
	\end{tabular}
	\caption{Registro delle modifiche.}
	\label{table:Registro delle modifiche}
\end{table}
\newpage
\tableofcontents
\clearpage
\newpage
\justifying
\section{Contenuti del Verbale}
\subsection{Informazioni sulla riunione}
\begin{itemize}
	\setlength\itemsep{0em}
	\item\textbf{Luogo:} Chiamata Google Meet;
	\item\textbf{Ora di inizio:} 15:00;
	\item\textbf{Ora di fine:}  16:00.
\end{itemize}
\begin{table}[ht!]
	\begin{minipage}[t]{0.5\linewidth}
		\centering
		\begin{tabular}{p{3cm} p{3cm}}
			\toprule
			\textbf{Partecipante} & \textbf{Durata presenza} \\
			\midrule
			Guglielmo Barison & 1.0 h \\
			Linda Barbiero &  1.0 h \\
			Pietro Busato & 1.0 h \\
			Oscar Konieczny & 1.0 h \\
			Davide Donanzan & 1.0 h \\
			Veronica Tecchiati & 1.0 h \\
			\bottomrule
		\end{tabular}
		\caption{Partecipanti NaN1fy.}
		\label{table:Partecipanti NaN1fy}
	\end{minipage} 
	\begin{minipage}[t]{0.5\linewidth} % -- COMMENTA/DECOMMENTA DA QUI
		\centering
		\begin{tabular}{p{3cm} p{3cm}}
			\toprule
			\textbf{Partecipante} & \textbf{Durata presenza} \\
			\midrule
			Andrea Dorigo & 1.0 h \\
			Fabio Pallaro &  1.0 h \\
			Daniele Zorzi & 0.0 h \\
			\bottomrule
		\end{tabular}
		\caption{Partecipanti SyncLab.}
		\label{table:Partecipanti SyncLab}
	\end{minipage} % -- A QUI PER TOGLIERE AGGIUNGERE
\end{table}

\subsection{Ordine del giorno}
\begin{itemize}
	\setlength\itemsep{0em}
	\item Revisione dello stato di avanzamento degli obbiettivi del secondo Sprint;
	\item Discussione sui problemi riscontrati nella configurazione di Grafana.
\end{itemize}
\subsection{Sintesi dell'incontro}
L'incontro con i membri della Proponente ha avuto come oggetto principale la disamina dei risultati ottenuti nel corso del secondo Sprint, con particolare attenzione all'implementazione e all'integrazione di Grafana con il prodotto finora ottenuto. Si è inoltre discusso di alcuni possibili miglioramenti nel contesto di stesura del codice ed esposizione dello stesso. Si è infine ampiamente discusso di un problema riscontrato dal team nella procedura di configurazione di Grafana, che non ha comunque impedito il raggiungimento degli obiettivi prefissati per la fine dello Sprint.

\subsubsection{Avanzamenti nel codice}
Gli obiettivi concordati erano l'implementazione di \textbf{Clickhouse nel container Docker}, di uno \textbf{stream dati da Kafka a Clickhouse} e una prima \textbf{visualizzazione dei dati di Clickhouse su dashboard Grafana}. Il team ha realizzato nel Docker precedentemente costruito il corretto inserimento di Clickhouse e ha verificato l'effettiva trasmissione dati da Kafka a Clickhouse; sono state riscontrate dalla Proponente alcune minori inconvenienze, velocemente risolvibili, quali l'erroneo fuso orario nei dati generati ed una dipendenza superflua di Clickhouse dallo stato di Kafka. Il team è infine riuscito nell'implementazione di una prima visualizzazione dei dati tramite Grafana, pur riscontrando alcuni problemi durante la fase di configurazione. La Proponente ha consigliato, per evitare confusione durante le esposizioni, di eseguire delle pull request per poter mostrare i risultati ottenuti senza cambiare costantemente da un branch all'altro della repository. La Proponente ha infine dichiarato che il prodotto mostrato era conforme alle aspettative per questo Sprint.

\subsubsection{Obiettivi dello Sprint}
La Proponente ha richiesto, per lo Sprint successivo, di provvedere alla realizzazione di una dashboard di Grafana più avanzata, con visualizzazioni di dati complessi o elaborati, e di query che coinvolgano più tabelle. Viene chiesto inoltre di generare nuovi dati, possibilmente di tipo diverso rispetto a quanto già implementato, visualizzabili anch'essi tramite Grafana. La fine del terzo Sprint, e di conseguenza la data del prossimo SAL, è stata fissata per venerdì 17 maggio. 

\subsection{Conclusioni}
\begin{itemize}
	\setlength\itemsep{0em}
	\item Viene fissato per il 2024-05-17 alle ore 15:00 il prossimo SAL;
	% consigliata la forma \textit{Viene adottato} quando viene adottato un certo modo di fare/strumento
	% per nomi di aziene e capitolati usare \texttt{}, e.g \texttt{Easy meal} \texttt{C6} 
\end{itemize}

\section{Attività da svolgere}
È di primaria importanza stabilire subito la divisione dei compiti assegnati per il prossimo
Sprint, oltre che delle stesure dei vari documenti per l'RTB. Per questo motivo, è indetta una riunione interna per definire i dettagli organizzativi del terzo Sprint. L'incontro è fissato per il giorno stesso, a cavallo tra la fine del SAL e l'inizio del ``Diario di Bordo" tenuto dal prof. Vardanega.
% commentare per togliere la firma
\signatureline{Padova, 2024-05-09}
\end{document}
